% !TeX root=../main.tex
% مشخصات فارسی رساله دکتری
%%%%%%%%%%%%%%%%%%%%%%%%%%%%%%%%%%%%
\university{دانشگاه تهران}
\college{دانشکدگان علوم و فناوری‌های میان‌رشته‌ای}
\faculty{دانشکدگان علوم و فناوری‌های میان‌رشته‌ای}
% گروه آموزشی خود را وارد کنید (در صورت نیاز)
\department{گروه علوم داده}
% رشته تحصیلی خود را وارد کنید
\subject{مهندسی شبکه‌های کامپیوتری}
% گرایش خود را وارد کنید
\field{؟؟}
% عنوان رساله
\title{بازسازی کور مشترک کدهای توربو با بازیابی همزمان کدگذارهای RSC و درهم‌ریزهای داخلی و خارجی}
% عنوان کوتاه برای سربرگ صفحات (تا با عنوان فصل تداخل نکند)
\runningtitle{بازسازی کور مشترک کدهای توربو}
% نام استاد(ان) راهنما
\firstsupervisor{دکتر مهدی تیموری}
\firstsupervisorrank{دانشیار}
%\secondsupervisor{دکتر راهنمای دوم}
%\secondsupervisorrank{استادیار}
% نام استاد(ان) مشاور؛ در صورت نداشتن مشاور، خطوط زیر را غیرفعال کنید.
\firstadvisor{دکتر کاکائی}
\firstadvisorrank{استادیار}
%\secondadvisor{دکتر مشاور دوم}
% نام داوران داخلی و خارجی
\internaljudge{دکتر داور داخلی}
\internaljudgerank{دانشیار}
\externaljudge{دکتر داور خارجی}
\externaljudgerank{دانشیار}
\externaljudgeuniversity{دانشگاه داور خارجی}
% نماینده کمیته تحصیلات تکمیلی
\graduatedeputy{دکتر نماینده}
\graduatedeputyrank{دانشیار}
% نام دانشجو
\name{شهریار}
\surname{توقی}
\studentID{۰۰۰۰۰۰۰۰۰}
\thesisdate{شهریور ۱۴۰۵}

%%%%%%%%%%%%%%%%%%%%%%%%%%%%%%%%%%%%%%%%%%%%%%%%%%%%
\dedication
{
{\Large تقدیم به:}\\
\begin{flushleft}{
	\huge
	\hspace{1em}\\
}
\end{flushleft}
}

\acknowledgement{
سپاس خداوندگار حکیم را که با لطف بی‌کران خود، آدمی را به زیور عقل آراست.

در آغاز وظیفهٔ خود می‌دانم از زحمات بی‌دریغ استاد(ان) راهنما و مشاور صمیمانه تشکر و قدردانی کنم که در طول انجام این رساله همواره راهنما و مشوق اینجانب بودند.

از خانواده و دوستانی که در این مسیر همراهی کردند نیز سپاسگزارم.
}

%%%%%%%%%%%%%%%%%%%%%%%%%%%%%%%%%%%%
\fa-abstract{
شناسایی و بازسازی کور ساختار کدهای کانال، به‌ویژه کدهای توربو، در مخابرات غیرهمکارانه یک مسئلهٔ کلیدی است. در این رساله، مسئلهٔ بازسازی مشترک پلی‌نوم‌های کدگذارهای RSC، درهم‌ریز داخلی و درهم‌ریز خارجی با درنظرگرفتن خاتمهٔ ترلیس و نویز صورت‌بندی و حل می‌شود. رویکرد پیشنهادی بر قیود جبری کدگذار، آزمون آماری قیود تحت نویز و جست‌وجوی درختی مقید استوار است تا پارامترهای مجهول از روی مشاهدات سخت (و در تعمیم‌ها، نرم) بازیابی شوند. نوآوری‌های اصلی شامل صورت‌بندی جامع مسئله، بازسازی همزمان مؤلفه‌های کد توربو، مدل‌سازی ابهام ساختاری و تحلیل رابطهٔ نویز با پیچیدگی محاسباتی است. نتایج شبیه‌سازی و تحلیل‌های نظری نشان می‌دهند روش پیشنهادی در شرایط نویزی نیز قادر به بازسازی ساختار مجاز کد است؛ جزئیات عددی پس از تکمیل ارزیابی‌ها در چکیده نهایی خواهد آمد.
}
\keywords{بازسازی کور، کد توربو، کدگذار RSC، درهم‌ریز، مخابرات غیرهمکارانه، شناسایی کور کد کانال}

\graphicalabstract{
\begin{center}
\textit{جای تصویر چکیده گرافیکی: معماری کلی رساله و چارچوب روش پیشنهادی بازسازی کور مشترک.}
\end{center}
}

\similarityresult{
\begin{center}
\textit{جای تصویر نتیجه سامانه تشابه‌یابی (همانندجویی).}
\end{center}
}

\fundingtext{
این رساله تنها با حمایت مالی دانشگاه تهران انجام شده است و حامی مالی دیگری نداشته است.
}
